\uuid{Dr0L}
\exo7id{3664}
\titre{exo7 3664}
\auteur{quercia}
\organisation{exo7}
\datecreate{2010-03-11}
\isIndication{false}
\isCorrection{true}
\chapitre{Espace euclidien, espace normé}
\sousChapitre{Produit scalaire, norme}
\cpge{mpsi-mp2i.espaces-prehilbertiens.a, pcsi.espaces-prehilbertiens.a, ptsi.geometrie-plan.b, mpsi-mp2i.espaces-prehilbertiens.b, pcsi.espaces-prehilbertiens.b, ptsi.geometrie-espace.b, mp-mpi.topologie-espaces-normes.a, pc.espaces-vectoriels-normes.a, psi.espaces-vectoriels-normes.a, psi.espaces-prehilbertiens.a, pt.structure-prehilbertienne.a, bcpst2.produit-scalaire.a}
\module{Algèbre}
\niveau{L2}
\difficulte{}

\contenu{
\texte{
Soit $E$ un espace euclidien.
}
\begin{enumerate}
    \item \question{Pour $\vec x \in E\setminus\{\vec 0\}$, on pose
    $f(\vec x) = \frac{\vec x}{\|\vec x\,\|^2}$.
    Montrer que : $\forall\ \vec x,\vec y \in E \setminus\{\vec 0\}$,
    $\|f(\vec x) - f(\vec y)\| = \frac{\|\vec x - \vec y\,\|}{\|\vec x\,\|\,\|\vec y\,\|}$.}
    \item \question{Soient $\vec a,\vec b,\vec c,\vec d \in E$.
    Montrer que $\|\vec a-\vec c\,\|\,\|\vec b-\vec d\,\| \le
         \|\vec a-\vec b\,\|\,\|\vec c-\vec d\,\| +
         \|\vec b-\vec c\,\|\,\|\vec a-\vec d\,\|$.


Indication : se ramener au cas $\vec a = \vec 0$ et utiliser l'application $f$.}
\reponse{
\'Elever au carré.
}
\end{enumerate}
}
